\documentclass[11pt]{letter}
\usepackage[margin=2.5cm]{geometry}
\usepackage{mathpazo}
\usepackage{hyperref}
\hypersetup{colorlinks=true,urlcolor=blue!70!black}

\signature{Theodor Spiro\\Independent Researcher, Paris, France\\theospirin@gmail.com}
\address{Theodor Spiro\\Paris, France\\theospirin@gmail.com}

\begin{document}
\begin{letter}{Editorial Board\\eLife Sciences Publications}

\opening{Dear Editors,}

I am submitting the manuscript ``Spectral Exponents of the Twelve-Lead ECG Reveal the Anatomy of Cardiac Conduction Disorders and a Bifurcation Between Aging and Disease'' for consideration as a Research Article at \textit{eLife}.

\textbf{Summary.} I apply IRASA---a method from computational neuroscience---to extract the spectral exponent $\beta$ from the aperiodic component of 412,730 twelve-lead ECGs across three continents. The central finding is a cross-population invariant: despite factor-of-two variation in absolute $\beta$-values across equipment, the twelve-lead $\beta$-vector distinguishes complete left from complete right bundle branch block with AUC = 0.982 on the German discovery cohort, AUC = 0.982 on a Chinese replication cohort, and AUC = 0.979 on a Brazilian replication cohort. This geometry is determined by the anatomy of the conduction system and is therefore invariant across populations.

\textbf{Why eLife.} I believe this work aligns with \textit{eLife}'s mission in three ways:

\begin{enumerate}
\item \textbf{Rigorous replication over novelty.} The core claim rests on identical replication across three independent datasets spanning three continents, recording systems, and populations---rather than on a single-dataset discovery. The analysis also includes an honest null result: $\beta$ does not independently predict mortality (adjusted HR = 1.02, $p = 0.83$), which I report because it calibrates what $\beta$ is (a diagnostic marker of conduction anatomy) and what it is not (a prognostic biomarker).

\item \textbf{Interpretability over performance.} The 12-lead $\beta$-vector provides an interpretable, physics-based characterization of conduction system integrity---each subtype produces a spatial fingerprint that maps directly onto lesion anatomy. This complements rather than competes with deep learning approaches that achieve marginal AUC improvements at the cost of interpretability.

\item \textbf{Open science.} All three datasets are publicly available. The analysis code is fully open: a Kaggle notebook for reproducibility (\url{https://www.kaggle.com/code/theodorspiro/the-heartbeat-at-the-edge-of-chaos}) and a GitHub repository (\url{https://github.com/mool32/ecg-spectral-exponents}). A preprint is available on bioRxiv [DOI to be added].
\end{enumerate}

\textbf{Calibration of claims.} The manuscript explicitly separates replicated findings (diagnostic landscape, spatial fingerprints, CLBBB vs.\ CRBBB AUC) from single-dataset observations (aging bifurcation, subclinical detection) and null results (mortality, absolute $\beta$ equipment-dependence). The aging trajectory did not replicate externally and is presented as a PTB-XL observation requiring independent validation.

\textbf{Suggested reviewers:}
\begin{itemize}
\item Bradley Voytek (UC San Diego) --- expertise in $1/f$ neural dynamics and IRASA methodology
\item Ary L. Goldberger (Harvard / Beth Israel) --- originator of the ``loss of complexity'' framework in cardiac physiology
\item Nils Strodthoff (Fraunhofer HHI) --- creator of the PTB-XL benchmark and deep learning ECG analysis
\item Antonio H. Ribeiro (Uppsala University) --- creator of the CODE-15\% dataset and DNN-based ECG age prediction
\item Hagen Schmidt (Radboud University) --- recent work on IRASA-based ECG spectral slopes and aging
\end{itemize}

Thank you for considering this manuscript.

\closing{Sincerely,}

\end{letter}
\end{document}
